% f03_claim1 variant d -- the negative control: a compliant reduction beside an
% order-violating one, at the same N, on the same bytes.
% Data: computed from src/sprs_core.py's own leaf generator (make_leaf_val,
%       alu_op=0b111) over the canonical heap-indexed tree. The canonical root
%       reproduces tables/invariance.tex row N=128 EXACTLY (0x415F7E8FF9581062),
%       which validates the model; left-to-right sequential summation over the
%       identical leaves then yields 0x415F7E8FF9581063 -- one ulp, one hex digit.
% Strategy: variants a-c argue that the root is INVARIANT across fabrics. This one
%       argues the converse half, which is what makes the claim non-vacuous: an
%       order that is not the canonical one gives a DIFFERENT answer. Without this
%       the invariance result could be read as measuring nothing.
%       (Both roots were re-derived independently before publication: the
%       canonical heap-indexed evaluation of the N=128 leaf vector reproduces
%       0x415F7E8FF9581062 exactly, and left-to-right summation of the same
%       leaves gives 0x415F7E8FF9581063.)
%       The leaf row is schematic -- eight leaves are drawn, while both roots
%       are computed over the full N=128 leaf vector.
% Colour: the compliant canonical tree is spBlue and its published root is the
% confirmed invariant (spGreen); the order-violating left-to-right chain and its
% one-ulp-different root are spVerm. Tree shape (log depth against linear chain)
% and the differing final hex digit carry the contrast without colour.
% Target: IEEE double column, 7.16in = 18.2cm. The body is drawn 9.2cm wide and
%       scaled 1.7x with transform shape, so it occupies 15.6cm.
\begin{tikzpicture}[
  scale=1.7, transform shape,
  font=\footnotesize,
  lf/.style={draw=spSky!80!spInk, rounded corners=1pt, minimum width=4.4mm,
             minimum height=3.6mm, inner sep=0pt, font=\tiny, fill=spSkyF, text=spInk},
  op/.style={draw=spBlue, circle, inner sep=0pt, minimum size=3.4mm, font=\tiny,
             fill=spBlueF, text=spInk},
  opv/.style={draw=spVerm, circle, inner sep=0pt, minimum size=3.4mm, font=\tiny,
              fill=spVermF, text=spInk},
  hdr/.style={font=\scriptsize\bfseries},
  arw/.style={-{Stealth[length=1.5mm]}, spBlue!75, thin},
  arv/.style={-{Stealth[length=1.5mm]}, spVerm!80, thin, densely dashed},
]
% ---------------- compliant: canonical tree ----------------
\node[hdr, anchor=west, text=spBlue] at (0,2.35) {compliant --- $\FBT(N)$};
\node[op] (r1) at (1.95,1.62) {$+$};
\node[op] (l1) at (0.85,1.02) {$+$};
\node[op] (l2) at (3.05,1.02) {$+$};
\node[op] (m1) at (0.30,0.42) {$+$};
\node[op] (m2) at (1.40,0.42) {$+$};
\node[op] (m3) at (2.50,0.42) {$+$};
\node[op] (m4) at (3.60,0.42) {$+$};
\foreach \i/\x in {0/0.05, 1/0.55, 2/1.15, 3/1.65, 4/2.25, 5/2.75, 6/3.35, 7/3.85}
  \node[lf] (f\i) at (\x,-0.18) {$L_{\i}$};
\foreach \a/\b in {f0/m1,f1/m1,f2/m2,f3/m2,f4/m3,f5/m3,f6/m4,f7/m4,
                   m1/l1,m2/l1,m3/l2,m4/l2,l1/r1,l2/r1} \draw[arw] (\a) -- (\b);
\node[font=\tiny, anchor=west, text=spBlue] at (2.25,1.90) {depth $\lceil\log_2 N\rceil$};

% ---------------- violating: left-to-right chain ----------------
\node[hdr, anchor=west, text=spVerm] at (4.9,2.35) {order-violating --- left to right};
\foreach \i/\x in {0/5.05, 1/5.55, 2/6.05, 3/6.55, 4/7.05, 5/7.55, 6/8.05, 7/8.55}
  \node[lf] (g\i) at (\x,-0.18) {$L_{\i}$};
\foreach \i/\x in {1/5.30, 2/5.80, 3/6.30, 4/6.80, 5/7.30, 6/7.80, 7/8.30}
  \node[opv] (c\i) at (\x,{0.30+0.26*\i}) {$+$};
\draw[arv] (g0) -- (c1); \draw[arv] (g1) -- (c1);
\foreach \i [evaluate=\i as \j using int(\i+1)] in {1,...,6} {
  \draw[arv] (c\i) -- (c\j); \draw[arv] (g\j) -- (c\j); }
\node[font=\tiny, anchor=north west, text=spVerm] at (4.95,2.10) {depth $N-1$};

% ---------------- the roots ----------------
\node[anchor=north west, font=\tiny\ttfamily, inner sep=1.5pt, text=spGreen]
  at (-0.05,-0.62) {0x415F7E8FF958106\textbf{2}};
\node[anchor=north west, font=\tiny\ttfamily, inner sep=1.5pt, text=spVerm]
  at (4.95,-0.62) {0x415F7E8FF958106\textbf{3}};
\node[anchor=north west, font=\tiny, text=spGreen, inner sep=1.5pt] at (-0.05,-1.02)
  {published root, $N{=}128$};
\node[anchor=north west, font=\tiny, text=spVerm, inner sep=1.5pt] at (4.95,-1.02)
  {same leaves, one ulp apart};
\draw[spSky!80!spInk, thin, <->] (4.15,-0.80) -- (4.90,-0.80);
\node[font=\tiny, anchor=south, inner sep=1.5pt, text=spSky!60!spInk] at (4.52,-0.66)
  {identical $L$};
% ---------------- what the contrast buys ----------------
\node[anchor=north, align=center, font=\tiny, text width=9.0cm, text=spInk] at (4.3,-1.42)
  {The left root is the one every fabric returns: over the whole campaign
   (\WallTBs{} settled simulations, \InvNonSettling{} non-settling) no
   configuration produced a finite root differing from $\Rcan(N,L)$, and at
   $N{=}128$ alone it survives six $(G,\text{topology})$ points over 371
   passing hardware configurations. The right root shows what that would have
   looked like had the order not been fixed.};
\end{tikzpicture}
